\documentclass[11pt]{article}

\usepackage[T1]{fontenc}
\usepackage[utf8]{inputenc}
\usepackage{amsmath,amssymb,amsthm,mathtools}
\usepackage[margin=1.12in]{geometry}
\usepackage{microtype}
\usepackage{booktabs}
\usepackage{xcolor}
\usepackage{hyperref}

\hypersetup{
  colorlinks=true,
  linkcolor=blue,
  citecolor=blue,
  urlcolor=blue,
  pdftitle={The First Multiset-Permutation Instability for the Consecutive Pattern 21435},
  pdfauthor={Xiaojian Zeng}
}

\newtheorem{theorem}{Theorem}
\newtheorem{lemma}[theorem]{Lemma}
\newtheorem{remark}[theorem]{Remark}

\newcommand{\W}{\mathcal W}
\newcommand{\one}{\mathbf 1}

\title{The First Multiset-Permutation Instability\\
for the Consecutive Pattern $21435$}
\author{Xiaojian Zeng\\
\small PBC School of Finance, Tsinghua University, Beijing 100190, China\\
\small \texttt{xiaojianzeng@foxmail.com}}
\date{}

\begin{document}
\maketitle

\begin{abstract}
For a positive multiplicity vector $\alpha=(\alpha_1,\ldots,\alpha_m)$,
let $F_\alpha(t)$ enumerate words of content $\alpha$ by occurrences of the
consecutive pattern $21435$.  We prove that $F_\alpha$ is invariant under
permuting the coordinates of $\alpha$ whenever the total length is at most
eight.  At length nine, the unique unstable unordered content is
$\{2,1^7\}$.  If the repeated letter has rank $1,2,3,$ or $4$, its polynomial
is
\[
  176026+5368t+46t^2,
\]
whereas ranks $5,6,7,$ and $8$ give
\[
  176025+5370t+45t^2.
\]
The proof classifies the four possible overlaps, evaluates the first two
binomial occurrence moments, and uses binomial inversion.  Exact exhaustive
enumeration through length nine is supplied only as a reproducibility check.
This result concerns one boundary case and does not resolve the general
classification of stable consecutive patterns.
\end{abstract}

\section{Context and statement}

The cluster method organizes consecutive occurrences by their overlaps
\cite{GouldenJackson1979}; in the permutation setting, overlap constraints can
be encoded by posets and their linear extensions
\cite{ElizaldeNoy2012}.  Generating functions for occurrences of consecutive
patterns in permutations and words have also been studied through symmetric
functions \cite{MendesRemmel2006}.  For multiset permutations, an earlier
symmetry result for classical pattern avoidance is due to Savage and Wilf
\cite{SavageWilf2006}.

Chen, Fang, and Kitaev recently formalized \emph{stability}: the full
occurrence distribution of a pattern is stable when it is unchanged by every
permutation of the multiplicities \cite{ChenFangKitaev2026}.  They proved a
broad sufficient criterion for instability.  Their Remark~4.13 singles out
$\underline{21435}$ because its minimal standard extension is the
all-distinct word $2143657$, so that criterion does not decide this pattern.
They conjectured that a consecutive pattern is stable exactly when it is
monotone.  The theorem below settles only the exact finite boundary for this
one pattern.  In particular, it does not settle their general conjecture.

Let $\alpha=(\alpha_1,\ldots,\alpha_m)\in\mathbb Z_{>0}^m$, put
$N=\sum_j\alpha_j$, and let $\W_\alpha$ be the words in which the letter $j$
appears $\alpha_j$ times.  For $w=w_1\cdots w_N$, define
\[
 c(w)=\#\{i: w_{i+1}<w_i<w_{i+3}<w_{i+2}<w_{i+4}\}
\]
and
\[
 F_\alpha(t)=\sum_{w\in\W_\alpha}t^{c(w)}.
\]
We call the unordered multiplicity type of $\alpha$ stable if
$F_\alpha=F_{\sigma\alpha}$ for every permutation $\sigma$ of the
coordinates.

\begin{theorem}\label{thm:main}
The consecutive pattern $21435$ has the following exact classification.
\begin{enumerate}
\item Every positive multiplicity vector with $N\leq 8$ has stable type.
\item At $N=9$, the only unstable unordered type is $\{2,1^7\}$.
\item For that type, let $\alpha^{(r)}$ denote the content on eight ordered
  letters in which the letter of rank $r$ occurs twice.  Then
  \[
  F_{\alpha^{(r)}}(t)=
  \begin{cases}
  176026+5368t+46t^2,&1\leq r\leq4,\\
  176025+5370t+45t^2,&5\leq r\leq8.
  \end{cases}
  \]
  No word of either content has three or more occurrences.
\end{enumerate}
\end{theorem}

The elementary symmetric functions entering the proof below are not an
asymptotic device: they are finite certificates that the marked-occurrence
counts are invariant under coordinate permutations.  For a distinct
coefficient-recovery setting in which elementary symmetric functions play an
algorithmic role, see \cite{Zeng2026Bezout}; we use no complexity result from
that work.

\section{The overlap posets}

An occurrence beginning at position $i$ imposes the strict chain
\begin{equation}\label{eq:one-occ}
 w_{i+1}<w_i<w_{i+3}<w_{i+2}<w_{i+4}.
\end{equation}
For two starts at distance $s\in\{1,2,3,4\}$, superimpose the two translated
copies of \eqref{eq:one-occ}.  This gives the complete overlap information
needed through length nine.

\begin{lemma}\label{lem:overlaps}
Two occurrences at start separation $1$ or $3$ are incompatible.  At
separation $2$, their relations form the seven-element chain
\begin{equation}\label{eq:shift2}
 w_{i+1}<w_i<w_{i+3}<w_{i+2}<w_{i+5}<w_{i+4}<w_{i+6}.
\end{equation}
At separation $4$, write
\[
\begin{aligned}
 a_1&=w_{i+1},&a_2&=w_i,&a_3&=w_{i+3},&a_4&=w_{i+2},\\
 c&=w_{i+4},&b&=w_{i+5},&d_1&=w_{i+7},&d_2&=w_{i+6},
 &d_3&=w_{i+8}.
\end{aligned}
\]
The resulting nine-element poset is
\begin{equation}\label{eq:shift4}
 a_1<a_2<a_3<a_4<c<d_1<d_2<d_3,
 \qquad b<c,
\end{equation}
and its only incomparable pairs are $\{b,a_j\}$ for $1\leq j\leq4$.
\end{lemma}

\begin{proof}
Use zero-based local positions.  A copy of \eqref{eq:one-occ} has value order
$1<0<3<2<4$.  At shift $1$, the union contains the directed cycle
\[
 2<1<0<3<2,
\]
and at shift $3$ it contains $3<2<4<3$.  At shift $2$, taking transitive
closure gives the position order
\[
 1,0,3,2,5,4,6,
\]
which is \eqref{eq:shift2}.  At shift $4$, writing the two chains explicitly
gives \eqref{eq:shift4}.  The element $b$ lies below $c$ but has no relation
to the four-element chain below $c$; all other pairs are comparable.
\end{proof}

\begin{lemma}\label{lem:degree}
If $N\leq9$ and the content of $w$ has a repeated letter, then $c(w)\leq2$.
\end{lemma}

\begin{proof}
Among words of length at most nine, two distinct occurrence starts differ by
at most four.  Lemma~\ref{lem:overlaps} shows that all pairwise differences
among three starts would have to lie in $\{2,4\}$.  The only possible
three-start set is therefore $\{i,i+2,i+4\}$, and it requires $N=9$.
Superimposing its three occurrence chains gives the strict total order
\[
 w_{i+1}<w_i<w_{i+3}<w_{i+2}<w_{i+5}<w_{i+4}
 <w_{i+7}<w_{i+6}<w_{i+8}.
\]
Thus its nine letters are distinct, contradicting the assumed repetition.
\end{proof}

\section{Binomial moments and stability}

Set
\[
 D(\alpha)=\prod_{j=1}^m\alpha_j!,
 \qquad
 e_k(\alpha)=\sum_{\substack{S\subseteq[m]\\|S|=k}}
                    \prod_{j\in S}\alpha_j,
\]
with $e_k=0$ when $m<k$.  Marking subsets of occurrences leads to the
binomial moments
\[
 B_k(\alpha)=\sum_{w\in\W_\alpha}\binom{c(w)}k.
\]

\begin{lemma}\label{lem:moments}
For $N\leq9$,
\begin{align}
 B_1(\alpha)
   &=\one_{N\geq5}\,(N-4)(N-5)!\frac{e_5(\alpha)}{D(\alpha)},
       \label{eq:B1}\\
 B_2(\alpha)
   &=\one_{N\geq7}\,(N-6)(N-7)!\frac{e_7(\alpha)}{D(\alpha)}
     +\one_{N=9}\,Q(\alpha),                         \label{eq:B2}
\end{align}
where $Q(\alpha)$ counts assignments of content $\alpha$ to the
shift-$4$ poset \eqref{eq:shift4}.  For repeated contents,
\[
 Q(\alpha)=0
 \quad\text{unless the unordered type is }\{2,1^7\}.
\]
For $\alpha^{(r)}$ of that type,
\begin{equation}\label{eq:Q}
 Q(\alpha^{(r)})=
 \begin{cases}1,&1\leq r\leq4,\\0,&5\leq r\leq8.
 \end{cases}
\end{equation}
\end{lemma}

\begin{proof}
Fix one start.  Choose the five distinct ranks used in
\eqref{eq:one-occ}.  A chosen five-set $S$ has a unique placement in the
chain, and the remaining positions can be filled in
\[
 \frac{(N-5)!}{\prod_j(\alpha_j-\one_{j\in S})!}
 =\frac{(N-5)!}{D(\alpha)}\prod_{j\in S}\alpha_j
\]
ways.  Sum over $S$ and over the $N-4$ starts to obtain \eqref{eq:B1}.

For a fixed pair of starts at separation $2$, the seven-chain
\eqref{eq:shift2} gives in the same way
\[
 (N-7)!\frac{e_7(\alpha)}{D(\alpha)}.
\]
There are $N-6$ such start pairs.  Separations $1$ and $3$ are impossible,
and within $N\leq9$ a separation-$4$ pair occurs only at $N=9$.  Its count is
$Q(\alpha)$, proving \eqref{eq:B2}.  Because these are counts of marked sets
of occurrences, a word having further unmarked occurrences creates no
overcounting issue.

It remains to analyze $Q$.  Equal labels in a strict poset assignment can
occupy only incomparable elements.  In \eqref{eq:shift4}, every incomparable
pair is $\{b,a_j\}$, and all such pairs share $b$.  Hence a non-injective
assignment of all nine positions can have exactly one repeated rank, of
multiplicity two.  This proves the vanishing assertion for every other
repeated content.

For $\alpha^{(r)}$, the equal positions must be $b$ and one $a_j$.  The lower
five positions $a_1,a_2,a_3,a_4,b$ consequently use four distinct ranks, all
below the four distinct ranks at $c,d_1,d_2,d_3$.  They must therefore be
ranks $1,2,3,4$.  The chain fixes their placement, and $b$ can equal the
specified repeated rank in exactly one way.  This is possible precisely when
$r\leq4$, giving \eqref{eq:Q}.
\end{proof}

\begin{proof}[Proof of Theorem~\ref{thm:main}]
The cases $N<5$ or $m<5$ are immediate: no occurrence is possible.  Now
suppose the content has a repeated letter.  By Lemma~\ref{lem:degree}, for
$N\leq9$ we may write
\[
 F_\alpha(t)=A_0+A_1t+A_2t^2.
\]
Since $|\W_\alpha|=N!/D(\alpha)$, binomial inversion is simply
\begin{equation}\label{eq:inversion}
 A_2=B_2,
 \qquad A_1=B_1-2B_2,
 \qquad A_0=\frac{N!}{D(\alpha)}-B_1+B_2.
\end{equation}

For $N\leq8$, every expression on the right of
\eqref{eq:B1}, \eqref{eq:B2}, and \eqref{eq:inversion} is symmetric in the
coordinates of $\alpha$.  Thus every repeated content has stable type.  An
all-distinct content has a singleton coordinate orbit, so it is stable as
well.

At $N=9$, the same argument applies to every repeated unordered type other
than $\{2,1^7\}$, because Lemma~\ref{lem:moments} makes its $Q$ term zero.
The all-distinct type again has a singleton coordinate orbit.  It remains to
evaluate the exceptional type.  For every $r$,
\[
 D(\alpha^{(r)})=2,
 \qquad |\W_{\alpha^{(r)}}|=\frac{9!}{2}=181440,
\]
and
\[
 e_5(\alpha^{(r)})=2\binom74+\binom75=91,
 \qquad
 e_7(\alpha^{(r)})=2\binom76+\binom77=15.
\]
Consequently
\[
 B_1=5\cdot4!\frac{91}{2}=5460,
 \qquad
 B_2=3\cdot2!\frac{15}{2}+Q(\alpha^{(r)})
     =45+Q(\alpha^{(r)}).
\]
Substitution of \eqref{eq:Q} into \eqref{eq:inversion} gives the two stated
polynomials.  They differ, so $\{2,1^7\}$ is unstable.  The degree assertion
is Lemma~\ref{lem:degree}.
\end{proof}

\section{Reproducibility and scope}

The accompanying standard-library Python program
\texttt{overlap\_moment\_check.py} reconstructs the transitive closures in
Lemma~\ref{lem:overlaps} and enumerates every word for all $511$ ordered
positive contents of total length at most nine.  It checks the two moment
formulas using exact integers over $7{,}685{,}705$ words, verifies that the
maximum occurrence count for repeated contents is two, and reproduces the
eight exceptional polynomials.  Its frozen plain-text output is included as
\texttt{overlap\_moment\_check.out}.  These computations corroborate the
argument; they are not used as a proof of the theorem.

The closest source comparison made here is to the cited version of
Chen--Fang--Kitaev.  That source states neither the length-nine classification
nor the two coefficient lists proved above.  This limited comparison is not a
claim of priority.  The broader conjecture that all non-monotone consecutive
patterns are unstable remains outside the present result.

\paragraph{Computational and AI assistance.}
A generative-AI system assisted with proof exposition, bibliography screening,
and drafting the exact verification program.  The displayed mathematical
argument was checked against the explicit overlap relations and integer moment
identities, and the deterministic computations were replayed during a separate
internal mathematical audit, which does not constitute peer review.  The
verification uses no floating-point arithmetic and no external mathematical
software.  Neither the AI output nor the successful computation is treated as
mathematical proof.

\bibliographystyle{alpha}
\bibliography{references}

\end{document}
